\begin{figure}[htbp]
\centering
\begin{tikzpicture}
\begin{axis}[bookbar, height=4.3cm, width=0.66\linewidth, bar width=15pt,
  ylabel={Share (\%)},
  title={The master's asymmetry},
  symbolic x coords={{tokens generated},{token-service GPUs}},
  xtick=data, ymin=0, ymax=115,
  xticklabel style={font=\scriptsize\color{inkSec}},
  nodes near coords, nodes near coords style={font=\scriptsize\color{inkSec}},
  legend style={at={(0.5,-0.30)}, anchor=north, legend columns=2},
]
\addplot[fill=sOne] coordinates {({tokens generated},10) ({token-service GPUs},88.7)};
\addlegendentry{master (1M-token context)}
\addplot[fill=sThree] coordinates {({tokens generated},90) ({token-service GPUs},11.3)};
\addlegendentry{sub-agents (128K context)}
\end{axis}
\end{tikzpicture}

\vspace{3mm}

\begin{tikzpicture}
\begin{axis}[bookaxis, height=5.4cm,
  ymode=log,
  ylabel={Capacity (BGPU, log scale)},
  title={The universal-reach ladder: committed against required},
  symbolic x coords={Singapore,Japan,{United States}},
  xtick=data, ymin=300, ymax=3e6,
  ytick={1000,10000,100000,1000000},
  yticklabels={$10^3$,$10^4$,$10^5$,$10^6$},
  xticklabel style={font=\small\color{inkSec}},
  ybar, bar width=7pt, xmajorgrids=false, enlarge x limits=0.36,
  every axis plot/.append style={draw=none},
  legend style={at={(0.5,-0.22)}, anchor=north, legend columns=2},
]
\addplot[fill=sOne] coordinates {(Singapore,1000) (Japan,39000) ({United States},112000)};
\addlegendentry{committed fleet}
\addplot[fill=sThree] coordinates {(Singapore,3200) (Japan,28300) ({United States},62300)};
\addlegendentry{rung 1: uniform baseline service}
\addplot[fill=sTwo] coordinates {(Singapore,8300) (Japan,72700) ({United States},160000)};
\addlegendentry{rung 2: full-population tiered service}
\addplot[fill=sFour] coordinates {(Singapore,57000) (Japan,503000) ({United States},1130000)};
\addlegendentry{modeled national requirement}
\end{axis}
\end{tikzpicture}

\vspace{2mm}

\begin{tikzpicture}
\begin{axis}[bookbar, height=4.2cm, width=0.72\linewidth, bar width=15pt,
  ylabel={BGPU per open-tier researcher},
  title={Provisioning against requirement, per researcher},
  symbolic x coords={Singapore,Japan,{United States},{requirement}},
  xtick=data, ymin=0, ymax=1.15,
  xticklabel style={font=\scriptsize\color{inkSec}},
  nodes near coords, nodes near coords style={font=\scriptsize\color{inkSec}},
  every node near coord/.append style={/pgf/number format/fixed,
     /pgf/number format/precision=2, /pgf/number format/fixed zerofill},
]
\addplot[fill=sOne] coordinates
  {(Singapore,0.02) (Japan,0.08) ({United States},0.10) ({requirement},1.00)};
\end{axis}
\end{tikzpicture}
\caption[The AI4S sizing results]{Three readings of the same framework~\cite{ai4s-framework}. \emph{Top}: ten
percent of tokens, generated at million-token context, consume nearly ninety percent of the token-service
capacity---the configuration's central structural fact, invariant across countries, and the reason master context
length is the layer's primary scenario axis. \emph{Middle}: committed national fleets against two priced service
rungs and the full modeled requirement, log scale. Japan's and the United States' commitments already exceed the
uniform-baseline rung but reach only $\sim$54\% and $\sim$70\% of the tiered rung; Singapore's clears neither.
\emph{Bottom}: the same result per researcher---announced infrastructures provision 0.02--0.10 BGPU against a
modeled requirement near 1.0, an underprovision of roughly one order of magnitude for Japan and the US and
closer to fifty-fold for Singapore. Japan is measurement-anchored [M/V]; Singapore and the US are
pre-measurement scenarios [T] whose behavioral parameters are transferred from the Japanese probe and are
labeled as such.}
\label{fig:ai4s-plots}
\end{figure}
